% page 476 du fac-simile (image p-475 du scan)
% bloc 154.9 x 246.3 mm ; marge gauche 8.0 mm
\pgc{-12.700mm}{0.000mm}{
\l{}
\l{\cel{3}468}
\l{\sou{protogina}. (\sou{biol.}) Floro che qua la pistilo matureskas ante}
\l{\cel{3}la stamini. (\sou{antonimo}\cel{1}: protandra). DEFIS.}
\l{}
\l{\sou{protokolo}. I. Raporto pri la rezolvi da kongresintaro, da}
\l{\cel{3}konferintaro diplomacala. - II. Raporto per qua persono}
\l{\cel{3}yurizita pri lo, o persono autoritatoza legale, konstatas}
\l{\cel{3}to quon lu agis, facis, vidis od audis. - III. Akto da}
\l{\cel{3}oficiero judiciala, qua konstatas delikto o kontravenco.}
\l{\cel{3}- IV. Raporto pri to quo eventis en+seanco o+sesiono di}
\l{\cel{3}asemblitaro. - DEFIRS.}
\l{}
\l{\sou{protoplasmo}. (\sou{biol.}) Substanco komplexa, pasable fluida e}
\l{\cel{3}diafana, ye qua konsistas la korpo di la celulo vivanta, e}
\l{\cel{3}qua kontenas parto diferenciita (\sou{nukleo}) e kelka organeti:}
\l{\cel{3}mikrosomi, leciti, vakuoli, e c. - DEFIRS.}
\l{}
\l{\sou{prototipo.}I.Tipo prima, segun qua onu exekutis la exempleri}
\l{\cel{3}cetera. - II. (\sou{metaf.}) Modelo supera ad la ceteri. - DEFIS}
\l{}
\l{\sou{protozoo}. (\sou{zool.)}\cel{2}Animaleto quan lua organismo elementala}
\l{\cel{3}lokizas ye la grado maxim infra en la hierarkio anima-}
\l{\cel{3}lala. - DEFIS.}
\l{}
\l{\sou{protuberanco}. Saliajo qua aspektas quale gibo, ye la surfaco}
\l{\cel{3}di korpo. - DEFIRS.}
\l{}
\l{\sou{proverbo}. Kurta precepto di sajeso praktikala, uzata da la}
\l{\cel{3}turbo. - DEFIS.}
\l{}
\l{\sou{providenco}. I. (\sou{religi}o)La sajo dea qua guvernas omna exis-}
\l{\cel{3}tajo. - II. (\sou{religio}). Deo egardata kom guvernanto di}
\l{\cel{3}omna existajo.\cel{2}III. (\sou{metaf}.)Helpo, susteno, konstanta.}
\l{\cel{3}EFIRS.}
\l{}
\l{\sou{provinco}. I. (en\cel{1}\sou{Roma, antique)}. Teritorio qua konquestesis}
\l{\cel{3}exter Italia, ed administrata da guberniestro. - II. Di-}
\l{\cel{3}viduro teritoriala di stato, administrata ye la nomo di}
\l{\cel{3}la suvereno da guverniestro partikulara. - III. Kontraste}
\l{\cel{3}kun la urbo metropola : la cetero di la lando. - DEFIRS.}
\l{}
\l{\sou{provizar}. (\sou{trans., ye ulo)}. Igar havar to quo esas necesa.}
\l{\cel{3}DEFIRS.}
\l{}
\l{\sou{provizora}. I. Qua supleas dum kelka tempo por la satisfaco}
\l{\cel{3}di bezono, varte di to quo esos definitiva. - II. Qua fun-}
\l{\cel{3}cionas en ofico, dum kelka tempo, sen havar la titulo qua}
\l{\cel{3}inheras la ofico. - DEFIRS.}
\l{}
\l{\sou{provokar.}\cel{1}(\sou{trans., ad).}\cel{2}Ecitar a lukto, pro agar atakeme.}
\l{\cel{3}DEFIRS.}
\l{}
\l{proxeneto. Mediacanto en la intrigi di aventuro amorala.}
\l{\cel{3}FIS.}
\l{}
\l{\cel{1}proxim. Poke distanta (tempale, spacale, gradale) de...EFIRS}
}
